% ==================== 第 8 章 项目管理 ====================
\section{项目管理与进度计划}

\subsection{团队分工}

项目分工按当前申报信息列示如下；提交前应由团队负责人复核成员姓名、
贡献比例和指导教师信息，确保与赛事报名表一致。

\begin{table}[h]
\centering
\caption{团队分工}
\label{tab:team}
\small
\begin{tabularx}{\textwidth}{l l l X}
\toprule
\textbf{成员} & \textbf{角色} & \textbf{贡献} & \textbf{工作量} \\
\midrule
李帅锋
  & 队长 / 算法 / 全栈
  & CP-SAT 求解器 / 评分卡 / 调课算法 / API 网关 / 前端 4 页 / CI
  & 约 60\% \\
张同学
  & 数据建模
  & 业务数据收集、演示数据构造、需求调研
  & 约 15\% \\
王同学
  & UI 设计
  & 设计稿、配色方案、用户访谈
  & 约 10\% \\
赵同学
  & 测试与文档
  & 单元测试 / 集成测试 / 用户手册撰写
  & 约 15\% \\
\bottomrule
\end{tabularx}
\end{table}

\subsection{进度计划}

\begin{table}[h]
\centering
\caption{第一阶段进度计划（2026.07.01--2026.07.12）}
\label{tab:timeline}
\small
\begin{tabularx}{\textwidth}{l l l X}
\toprule
\textbf{阶段} & \textbf{时间} & \textbf{里程碑} & \textbf{交付物} \\
\midrule
需求调研
  & 07.01--07.03
  & 梳理典型排课流程
  & 需求清单、对象模型 \\
数据建模
  & 07.04--07.05
  & 完成 5 类数据类
  & \code{models.py}、\code{demo.json} \\
硬约束求解
  & 07.06--07.07
  & 8 课次 0 冲突
  & \code{solver.py} v1 \\
软目标与多策略
  & 07.08--07.09
  & 5 套策略可比较
  & \code{solver.py} v2、调课算法 \\
REST API
  & 07.10
  & 6 个端点全通
  & \code{api.py}、CI 配置 \\
前端 4 页面
  & 07.11
  & 仪表盘 + 课表 + 调课 + 资源
  & \code{web/*.html/css/js} \\
测试与文档
  & 07.12
  & 19 测试通过
  & \code{tests/}、本报告 \\
\bottomrule
\end{tabularx}
\end{table}

\subsection{第二阶段规划}

\begin{table}[h]
\centering
\caption{第二阶段规划（2026.08--2026.12）}
\label{tab:phase2}
\small
\begin{tabularx}{\textwidth}{l l l X}
\toprule
\textbf{阶段} & \textbf{时间} & \textbf{目标} & \textbf{交付物} \\
\midrule
飞书集成准备
  & 08.01--08.15
  & 申请企业自建应用、获取 tenant\_access\_token
  & 飞书应用凭据、回调地址 \\
多维表格双向同步
  & 08.16--09.15
  & 业务数据可编辑、课表可回写
  & 多维表格 schema、同步服务 \\
消息卡片推送
  & 09.16--10.15
  & 方案确认后自动 @ 师生
  & 消息卡片模板、推送任务 \\
审批流程
  & 10.16--11.15
  & 调课需教务主任审批
  & 审批定义、Webhook 处理 \\
自然语言解析
  & 11.16--12.15
  & 飞书 AI 助手直接调用
  & Prompt 模板、调用 demo \\
总结与答辩
  & 12.16--12.31
  & 完成项目总结、准备答辩
  & 结题报告、演示视频 \\
\bottomrule
\end{tabularx}
\end{table}

\subsection{风险与对策}

\begin{table}[h]
\centering
\caption{主要风险与对策}
\label{tab:risk}
\small
\begin{tabularx}{\textwidth}{l c X X}
\toprule
\textbf{风险} & \textbf{概率} & \textbf{影响} & \textbf{对策} \\
\midrule
OR-Tools 性能不足
  & 中
  & 大规模场景超时
  & 1) 提前预筛选变量；2) 拆分大模型；3) 退而求其次用贪心 \\
飞书 API 配额不足
  & 中
  & 第二阶段功能受限
  & 申请高校合作计划免费额度 + 异步批处理 \\
教师/学生数据隐私
  & 中
  & 合规风险
  & 演示数据全虚构，商用前做数据脱敏 \\
团队成员变动
  & 低
  & 进度延迟
  & 文档与测试覆盖到位，新成员 1 周可接手 \\
业务需求变化
  & 中
  & 设计返工
  & 软目标权重表可灵活调整，硬约束通过数据驱动 \\
\bottomrule
\end{tabularx}
\end{table}

\subsection{本章小结}

本章给出了项目团队分工、第一阶段迭代进度、第二阶段
5 个月规划以及主要风险与对策。项目按计划顺利推进，
为第二阶段飞书生态集成打下了坚实基础。
